% Level 1, Buffer session B2 -- Free-Build and Remix Studio   (input by lesson-plans.tex)
\kcpart{B2}{Free-Build and Remix Studio}{Take something you made, change it on purpose, and make it yours.}
\lessoninfo{Concepts: remixing (I), all concepts so far (R/E), giving credit (R)}{Tools: Thonny + turtle (or Scratch, see note)}{CSTA: 1B-AP-11, 1B-AP-12, 1B-AP-14}

\begin{teachernote}[When to use B2]
Any time after week 8: after a missed session, before a holiday, between weeks 11 and 12 to finish projects, or when the group needs a change of pace. Before week 7, run the same plan in Scratch: students remix one of their own projects from weeks 1--5 (\menu{File} → \menu{Save to your computer} with a new name first).
\end{teachernote}

\begin{objectives}
  \item The student can choose an earlier program and change it in at least two deliberate ways (L1.G1, L1.G2; L1.WB2.1)
  \item The student can save a remix as a new file so the original stays safe (L1.G9; L1.WB2.2)
  \item The student can give credit for ideas and code they reused (L1.G8; L1.WB2.3)
  \item The student can find a bug a partner hid in a program (L1.G6, L1.G7; L1.WB2.4)
\end{objectives}

\begin{materials}
  \item One six-sided dice (or six folded slips in a cup, numbered 1--6). The remix menu is in the box \enquote{Remix dice} below --- show it on this page.
  \item On each Mac: Thonny open; Finder at \texttt{kids-coding/1-introduction/exercises/}; the folder \texttt{Documents/my-python}.
  \item Sticky notes for the gallery share. Prep time: about 5 minutes.
\end{materials}

\needspace{20\baselineskip}
\section{At a glance}
Times are a guide: free-build days can be flexible.
\begin{timeplan}
  \timeslot{0--8}{Warm-up}{Unplugged \enquote{Remix dice}: roll twists for a past exercise.}
  \timeslot{8--15}{Concept}{What a remix is; save as a new file; credit comment.}
  \timeslot{15--25}{Guided demo}{Remix the Square Flower live: random colors and a question.}
  \timeslot{25--48}{Hands-on build}{Pairs remix a past exercise with at least two twists; bug-swap at minute 40.}
  \timeslot{48--60}{Wrap-up}{Gallery share with \enquote{two stars and a wish}; credit check; reflection.}
\end{timeplan}

\section{Warm-up (0--8 min)}
\begin{unplugged}[Remix dice]
Each pair chooses a past exercise they liked (show the list: first square, polygon party, square flower, pick a color, how many stars, flower stamp, bubble art). Then they roll the dice twice; each number is a twist to add:
\begin{center}\small
\begin{tabular}{@{}cl@{\qquad}cl@{}}
\toprule
1 & Change colors (use the palette) & 4 & Ask a question with \texttt{input} \\
2 & Change sizes with a variable & 5 & Add a \texttt{random} surprise \\
3 & Add or change a loop & 6 & Make a \texttt{def} command \\
\bottomrule
\end{tabular}
\end{center}
Same number twice? Roll again or choose freely. Twists 5 and 6 need week 10 skills; before week 10, re-roll a 5 or 6. Pairs write their exercise and two twists on a sticky note: that is their plan.
\end{unplugged}

\section{Concept (8--15 min)}
\begin{teacherscript}
\enquote{A \textbf{remix} takes something that exists and changes it on purpose to make something new. Musicians do it, artists do it, and programmers do it all the time. Two rules: first, \textbf{save a copy} with a new name, so the original stays safe. Second, \textbf{give credit}: a comment at the top says where it came from.}
\end{teacherscript}
\begin{analogy}[A cake recipe with your own twist]
You take Grandma's chocolate cake recipe and add raspberries and a different icing. It's your new cake --- but when people love it, you say \enquote{it's Grandma's recipe with my twist}. That's a remix with credit.
\end{analogy}
\begin{safety}[Credit and remix]
The credit line goes at the very top: \texttt{\# Remixed from Exercise 8.3 (Square Flower) by Kai and Zoe}. Use nicknames. If you remix a friend's program, ask them first and name them. Taking someone's program and saying you made it alone is not fair --- even if you changed the colors.
\end{safety}

\section{Guided demo (15--25 min)}
\begin{workedexample}[Remix the Square Flower]
\begin{enumerate}[leftmargin=1.4em,itemsep=1pt]
  \item Open \texttt{week08/ex-square-flower/starter.py} (or a student's finished one). Before changing anything: \menu{File} → \menu{Save as…} → \texttt{Documents/my-python/square-flower-remix.py}. \enquote{Now the original is safe.}
  \item Add the credit comment at the top.
  \item Twist 4 (question): \texttt{petals = int(input("How many squares? (4 to 36) "))}. Change the loop to \texttt{for i in range(petals):} and the turn to \texttt{t.left(360 / petals)}. Run with 12.
  \item Twist 5 (random): \texttt{import random} at the top and \texttt{t.color(random.choice(PALETTE))} inside the loop. Run twice --- different colors each time.
\end{enumerate}
\end{workedexample}
\begin{pycode}[square-flower-remix.py (tested)]
# Square Flower REMIX
# Remixed from Exercise 8.3 (Square Flower) by Kai and Zoe
import random
import turtle

PALETTE = ["#E69F00", "#56B4E9", "#009E73", "#F0E442",
           "#0072B2", "#D55E00", "#CC79A7", "#000000"]

t = turtle.Turtle()
t.pensize(3)
t.speed(0)

petals = int(input("How many squares? (4 to 36) "))

for i in range(petals):
    t.color(random.choice(PALETTE))
    for side in range(4):
        t.forward(100)
        t.left(90)
    t.left(360 / petals)

turtle.done()
\end{pycode}

\section{Hands-on build (25--48 min)}
\begin{pairroles}[Driver and navigator]
Driver types; navigator holds the sticky-note plan and checks each twist off. Swap every 8 minutes. With 1 student: you are the navigator and roll the dice with them. With 3: rotate driver, navigator and \enquote{tester}. With 4: two pairs.
\end{pairroles}
\begin{bughunt}[Bug swap (minute 40)]
Each pair saves one more copy (\texttt{...-bug.py}) and hides \textbf{one} sneaky bug in it: a typo, a wrong angle, a line pulled out of a loop, or a missing \texttt{int}. Pairs swap Macs and have 3 minutes to find and fix it, then explain the bug to its makers. With 1 student, you hide a bug in their program and they hide one in yours. The rule: a bug must be fixable by changing one line.
\end{bughunt}
\begin{freebuild}[Your remix, your way]
After the two twists, anything goes: more twists, a mash-up of two exercises (flowers in the bubble art, stars in the polygon party), or work on the final project if it needs finishing before week 12.
\end{freebuild}

\section{Wrap-up (48--60 min)}
Gallery share: each pair runs their remix and says which exercise it came from and which twists they added. Others give \enquote{two stars and a wish} on sticky notes.
\begin{formativecheck}[Credit check]
Ask every pair to show the top line of their remix. Look for: a credit comment, a new file name, and at least two twists that work. Quick question: \enquote{Why did we use Save as before changing anything?} (So the original stays safe.)
\end{formativecheck}
\begin{reflection}
\begin{itemize}[leftmargin=1.2em,itemsep=1pt]
  \item Which twist changed your program the most?
  \item What would you remix next time, and how?
\end{itemize}
\end{reflection}

\section{Notes for the teacher}
\begin{commonmistakes}
  \item \textbf{Editing the original instead of a copy.} Fix: \menu{File} → \menu{Save as…} first, every time. If it already happened, nothing is lost: the exercise folder can be copied again from your Mac.
  \item \textbf{Twist breaks the program}: e.g. \texttt{PALETTE} or \texttt{import random} missing after copy-paste (\texttt{NameError}). Fix: read the last line of the error; add the missing line at the top.
  \item \textbf{Changing the number of sides but not the angle.} Use \texttt{360 / sides} so both always match.
\end{commonmistakes}
\begin{differentiation}
\supportpath Choose the first square or the polygon party and roll once; twist 1 or 2 is plenty. Sit as navigator.
\extensionpath Roll three times, or remix the final project's function into a whole new scene. Try making the remix react to a key: \texttt{turtle.onkey(...)} with \texttt{turtle.listen()}.
\end{differentiation}
\begin{stretchchallenge}[Remix of a remix]
Swap remixes with the other pair (with permission!), add one more twist, and add a second credit line: \texttt{\# Remix of Kai and Zoe's remix by Sam and Mei}.
\end{stretchchallenge}
\begin{teachernote}[If you are short on time]
Skip the dice: everyone adds random colors to one past exercise. Keep the credit line and the gallery share.
\end{teachernote}
